\chapter{Relevant Literature}

The term \textit{Privacy} by definition can be categorized in a broad manner, without explicitly ever encountering the actual depths of what it means to an individual. This complexity increases when the ideologies of culture is borrowed and intertwined in this context, to help understand the nuances of the environments it exists within. Such nuances have a lasting impact in the way people perceive the privacy of one's self, depending on the context that is being considered. Personal data privacy therefore, should ideally incite the strongest reactions of protection and safeguarding, since it involves dealing with information which is closest and dearest to the user with respect to their everyday lives. Past works \cite{zou_ive_nodate, mayer_now_nodate, karunakaran_data_nodate} however, show us that while the users display mixed awareness about the breach of their personal data privacy, more often than not they believed that the breach would not impact them in any reasonable manner. This elicits a notion of 'Privacy Paradox' \cite{acquisti_economics_2016}: people claiming that they are concerned about their privacy but often not taking measures to safeguard it. 

Regulations such as GDPR and CCPA were introduced to help the users by providing the necessary safeguards to take control over the provision, processing, and usage of their personal data by various organizations. While there was some initial skepticism towards the effectiveness of the individual claims and rights by the public \cite{strycharz_data_2020}, they were hopeful in terms of the implementation of such safeguards and maintained that it would eventually lead to more privacy-focused infrastructural developments \cite{mangini_empirical_2020}. Although the intention was to help the masses have an understanding over their control and rights, it has not translated into actual privacy practices. This left the users having a lack an understanding of the privacy disclosures due to the ambiguity on why their data is being collected \cite{chen_fighting_2021}, alternative ways for organizations to be notoriously hidden with their data practices \cite{oconnor_clear_2021}, disclosure of irrelevant information and overwhelming the users when they explicitly want practices to be transparent and fair. \cite{kyi_it_2024}. 

Schomakers et al. \cite{schomakers_internet_2019} compared the privacy perceptions of German users with those from the Brazilian and the US contexts, to understand if there any differences in their perspectives on the sensitivities of the different data items. They observed that although the ranking remained similar across the nations, there were differences in terms of the associated risks involved in disclosure of such information and how these could possibly be explained by factors such as habits, education level and privacy disposition. Similarly, Herbert et al. \cite{herbert_world_2023} conducted a large-scale longitudinal survey to understand if demographic factors played a role in the way that people held misconceptions about various security and privacy practices. They noticed that the country of residence played a key role in withholding such beliefs, while also suggesting the possibility of Internet literacy or technological readiness influencing the participants. Herbert et al. \cite{herbert_digital_nodate} also conducted another study using this same dataset, to understand the differences in the digital security \& privacy perceptions and practices of participants from WEIRD vs NON-WEIRD countries, focusing on their sensitivity of various data items, risks by possible attackers, previous experiences etc. They found that participants from non-WEIRD countries scored higher on all their measurement variables showcasing higher awareness, perception and exposure in terms of their security practices. One of their key findings was that the non-WEIRD participants relied heavily on their friends and family for digital security advice, which showcases a communal perspective towards these practices.

In an identical notion but with a qualitative perspective, Umbach et al \cite{umbach_your_2023} compared user perceptions on their security \& privacy practices of participants from Brazil, Philippines, India, Egypt and Nigeria. They found that one of the primary concerns of the participants revolved around their financial data and being wary of scammers in this regard, while also talking about their device-sharing methods obstructing with their individual right to privacy. This behavior was also noted across other contexts with Ahmed et al. \cite{ahmed_everyone_2019, ahmed_digital_2017} finding that while economic need is driving factor in device-sharing, it also stems from deep-rooted social and cultural practices in the Bangladeshi societies. They also found that existing power dynamics and patriarchal practices were also amplified due to this practice, thereby revealing its impact on the women's privacy autonomy in the society. Such dynamics and power structures in the context of personal data privacy was even more prevalent in the Arab Gulf, where Abhokhodair et al. \cite{abokhodair_privacy_2016, abokhodair_photo_2017} observed that decision making on Social media and Photo-sharing platforms involved not just the user's concerns of their safety and privacy but also an adherence to their religious and societal beliefs. In an almost similar environment, Naveed et al \cite{naveed_ask_2022} conducted interviews with Pakistani users from low-literacy, low-income backgrounds regarding their privacy choices and its implications on their personal data. They echoed similar sentiments of religious and societal beliefs of patriarchy playing a key role in especially impacting the women's management of privacy choices and control on their 'personal data'.

Kumaraguru et. al \cite{kumaraguru_privacy_2012} conducted an empirical study on the user perceptions of their privacy choices with a sample of 10,427 participants across India, and found that the participants were split in their opinions on what constituted as personal data and thereby more sensitive and private in nature. They also observed that the participants were considering financial data to be one of the most important forms of personal information, due to their awareness of the various financial scams. However, this work was conducted in 2012, making the results non-replicable given the growth of Internet penetration, awareness as well as the changes in the judicial aspects with respect to personal data privacy of the country since that time period. Recent works focused more on specific aspects of user privacy perceptions such as: Privacy practices of older house members and their families \cite{murthy_individually_2021}, Security and Privacy of UPI users \cite{mungara_security_nodate}, Geosocial Networking apps and the navigational challenges involved in them while safeguarding one's privacy \cite{pinch_someone_2022}, low-income users and their challenges to safeguard their privacy while navigating the identity infrastructure \cite{srinivasan_poverty_nodate}, low-income new Internet users and how they navigate through the realm of privacy \cite{popuri_i_nodate} and User's attitudes on AI sources \cite{kapania_because_2022}. 

Almost of these works dealt exclusively on specific aspects of user privacy perceptions or were conducted in specific contexts. Given the recency of the DPDPA's partial enforcement, it is therefore necessary to understand user privacy in the context of their personal data and how it is perceived. Priyadarsini et al. \cite{patil_web_2025} focusing primarily on the user's web privacy perceptions did touch this topic slightly, finding out that in their sample size, the participants were skeptical about sharing personal data like financial information and were more aware of sharing it on social-media applications than on e-commerce platforms. Additionally, Athar et al. \cite{10.1145/3779208.3785287} tried to understand the participant's interactions with cookie banners and how they perceived some of the laws of the DPDPA. They observed the possibility of 'privacy paradox' being present within their responses and an overall sentiment of distrust in the government's application of the Act. However, it is evident that while there has been extensive research done in the Indian ecosystem with regards to the user privacy perceptions, there has not been any work done focusing solely on user privacy of their personal data and its processing which is the primary focus of this Act. To the best of my knowledge, this is the first work which tries to understand the general privacy perceptions of the user with respect to their personal data along with their sensitivities of the various personal data items, helping us understand if there is a mismatch in the user comprehension aspect to what the Act is expected to cater. Minimizing this mismatch, could help in the progression of this Act to a wider community, thereby serving its purpose to the fullest.


